\chapter{Discussion}\label{ch:discussion}

\section{Main findings}
The reproduction succeeded. The survey prevalences of \PrevDep\% for probable depression and \PrevAnx\% for probable anxiety match the original study exactly, and the rise in women's diet diversity agrees with the panel authors' conclusion. The checks changed several of the interpretations.

The weaker mobility response to stringency in 2022 held up in first differences and without March 2020. The rise in diet diversity was the same in every wealth quintile within sampling error, so the poorer women did not fall further behind during the year, and the wealth gap that remains is the one that existed in 2019. A uniform upward shift is also what a change from in-person to telephone interviewing could produce. In the mental-health survey the associations are smaller than the odds ratios suggest. For a family death, the adjusted odds ratio of \OrFamDeathB{} corresponds to an adjusted prevalence ratio of \PrFamDeathB. The Chattogram and vaccination estimates depend partly on the flagged records, and the latent profiles are bands of one severity score. For prediction, the four learners could not be told apart, and the AUC depended almost entirely on which features were admitted.

The three datasets cannot be linked, so the national series serves only as context for the two surveys. What the two surveys have in common is modest. In both, characteristics that pre-date the pandemic distinguished respondents more than the measured pandemic exposures did. In the panel, the 2020 constraints were not associated with diet diversity once wealth and baseline diet were in the model (purchasing power \PanPPPp, market access \PanMKTp). In the mental-health survey, division, mental-health history and SES ranked above every COVID-19 variable in the tier-A importance.

\section{Comparison with previous studies}
My estimates agree with \textcite{islam2023} on prevalence and on the direction of the associations with student status, mental-health history, mindfulness, optimism and vaccination. The re-analysis qualifies that study in several respects. Regional prevalence does not follow a simple geographic pattern, and the low Chattogram figure depends in part on the flagged records. Sex is associated with anxiety but not with depression (\DepFemale\% of women and \DepMale\% of men met the threshold). The association with the respondent's own hospitalisation cannot be assessed from \MhHospN{} hospitalised respondents. The two-dose association weakens from an odds ratio of \OrTwoDose{} to \SensVaxExcl{} once the flagged records are set aside.

High symptom levels were also reported by other Bangladeshi surveys during the pandemic \parencite{islam2020,banna2022}, and global estimates place the largest increases in depression and anxiety where infection rates were high and mobility fell most \parencite{santomauro2021}. This survey has no pre-pandemic baseline and no probability sample, so it can neither confirm nor contradict an increase.

The household result contrasts with early warnings of a food-security crisis driven by lost income \parencite{laborde2020} and with evidence of rising food insecurity in Bangladesh immediately after the first lockdown \parencite{hamadani2020}. Timing may explain part of the difference, since the follow-up took place eight months after that lockdown, and survey mode may explain some of the rest. The panel authors themselves reported the rise in diet diversity \parencite{panel2021}. My addition is that the rise was flat across wealth quintiles and that the wealth gap is a pre-pandemic feature.

The learner comparison agrees with the systematic review by \textcite{christodoulou2019}, which found no performance benefit of machine learning over logistic regression for clinical prediction models. Calibration did differ between learners, and \textcite{vancalster2019} argue that it deserves as much attention as discrimination.

\section{Methodological lessons}
On the fixed hold-out split, the learner selected by cross-validation in tier B had the lowest test AUC of the four, and the cross-validated AUC exceeded the test AUC for every learner. Read alone, that looks like overfitting during tuning. The repeated splits show it is sampling variation, since untuned logistic regression has the same gap and the gap reverses sign in tier A. With \MlTestN{} test cases, a learner comparison needs repeated splits or an equivalent resampling scheme before any ranking is reported.

The ablation was more informative than the learner comparison. It showed that one scale carries the tier-B gain, that division carries a good part of tier A, and, through the PHQ-2 reference, how much AUC a feature earns just by overlapping with the outcome. Without that ceiling I would have read an AUC of \AblB{} from a short questionnaire as more useful than it is.

The flagged records were found by a mixture model, and the regressions gave no sign that \SuspectN{} rows were near-identical. A direct count of duplicated response blocks is the simpler audit, and I would now run it before any modelling.

Model form mattered in the other two datasets. In the national series the slopes are similar in levels and in first differences, but only the differenced model, whose residuals are close to uncorrelated, supports a test of the change between years. In the panel, the ANCOVA quintile coefficients read like larger gains for richer women while the change scores show equal gains, and a baseline coefficient of \PanBaseb{} explains why \parencite{lord1967}.

\section{Implications}
The wealth gap in women's diet diversity in these two districts is a standing condition. It appears in both years and under both interview modes, so programmes aimed at dietary adequacy among the poorest households have a case that does not depend on any estimate of the pandemic's effect.

Among the survey's respondents, depressive symptoms were concentrated in identifiable groups, namely people with a prior mental-health history, students and the bereaved. A representative survey would be needed to say how large these groups' excess is in the population, and nothing here shows what would reduce it. The data do indicate where such a survey, or a service with limited capacity, should look first.

\section{Limitations}
The most serious limitations concern design. The mental-health survey is an online convenience sample without weights, in which \MhDhaka\% of respondents live in Dhaka division and \MhStudents\% are students, so its prevalences describe the respondents and should not be read as national estimates. It is cross-sectional, its outcomes are screening thresholds on self-report scales, and none of its associations has a causal reading. In the panel, the change from in-person to telephone interviewing is fully confounded with the pandemic year, its direction of bias is unknown, and the panel covers two districts.

Data quality comes next. The screen catches one obvious pattern, identical blocks shared by five or more respondents, and other low-quality responses may remain. I do not know whether the flagged records are genuine, and the Chattogram and vaccination estimates depend on that. The PSS-4 and the pessimism sub-scale have low internal consistency, which weakens any result involving them.

The remaining limits are of scope. The prediction models were tested only on respondents from the same survey, and nothing is known about their performance in another sample, period or country. The three datasets are not linked, so no joint model across them is possible. The drift in Google's mobility baseline (Section~\ref{sec:retreat}) makes 2022 mobility levels partly an artefact of the reference period, a problem the weekly changes largely avoid.
